\section{Accountability: who is responsible when it is wrong?}
\label{sec:accountability}

A clinician must know where responsibility sits. The seventh question asks who
answers when something goes wrong, what is documented and by whom, and how
liability is allocated across the human and the system. A diffuse answer is no
answer: if everyone is responsible, no one is, and the clinician is left holding a
risk that was never theirs alone to carry. Accountability in this framework is
therefore drawn explicitly, with each role named and each obligation written down
before the trial begins.

The instrument that draws it is the standard operating procedure, the written,
signed, version-controlled document that fixes who does what, under which
thresholds, and in what order. Because the standard operating procedure is signed
by named parties, it converts an abstract division of labor into a record of who
agreed to what. Sitting alongside it is the audit trail, the append-only,
hash-chained record in which each action, gate result, and countersignature is
cryptographically linked, so that any action can be traced from outcome back to the
human or system responsible. The author's site documentation package specifies the
contents and custody of these records \cite{Kawchak2026PhysicalAIOncologyClinical},
and the regulatory-adaption lineage maps the roles onto existing frameworks: the
adaption of 21 CFR Part 312 places sponsor and investigator duties on a familiar
legal footing \cite{Kawchak2026Adaption21CFRPart}, while the adaption of the ICH
harmonised guideline aligns the same roles with international practice
\cite{Kawchak2026AdaptionICHHarmonisedGuideline}.

\autoref{tab:accountability} sets out, for each party, what that party is
responsible for and where the responsibility is documented. The clinician owns the
decision to accept, escalate, or override, recorded by countersignature in the
audit trail. The sponsor owns the protocol and oversight, the developer owns the
verified behavior of the system, the IRB owns the ethical approval and its
conditions, and the auditor owns independent confirmation that the record matches
what occurred. Figure 9 shows the same structure converging on the signed standard
operating procedure, with the auditor positioned to read across all of it.

This is what lets a clinician answer the question that matters when an action is
later challenged. If an auditor asked me tomorrow, could I show exactly what
happened and who was responsible? Yes, because the standard operating procedure
names the responsible party for each function and the audit trail binds each action
to that party. Accountability drawn this way also protects the clinician, because
it makes visible the obligations that belong to the sponsor and the developer
rather than letting them settle silently onto the person at the bedside.

\begin{cfwfig}
\node[cfact] (a) at (0,0) {Clinician}; 
\node[cftwo] (b) at (0,-1.3) {Sponsor};
\node[cftwo] (c) at (0,-2.6) {Developer}; 
\node[cfhope] (d) at (5.0,-1.4) {Signed SOP};
\node[cfone] (e) at (9.4,-1.4) {Auditor};
\draw[cfedge] (a.south east) -- (d.north west); 
\draw[cfedge] (b.east)       -- (d.west); 
\draw[cfedge] (c.north east) -- (d.south west); 
\draw[cfedge] (d.east)       -- (e.west);
\end{cfwfig}
\vspace{-0.7cm}
\cfwcaption{Figure 9. Where accountability sits: clinician, sponsor, and developer obligations converge on a signed standard
operating procedure that the auditor reads across.}

\needspace{9\baselineskip}\begin{table}[H]
\footnotesize
\begin{tabularx}{\textwidth}{@{}L{2.25cm} Y L{5.2cm}@{}}
\toprule
\textbf{Party} & \textbf{Responsible for} & \textbf{Documented in}\\
\midrule
Clinician & The decision to accept, escalate, or override an action & Countersignature in the audit trail\\
Sponsor & The protocol, oversight, and adverse-event response & Signed standard operating procedure\\
Developer & The verified behavior of the system and its gates & Verification record and gate results\\
IRB & Ethical approval and the conditions attached to it & App. letter, standard operating proc.\\
Auditor & Independent confirmation that record matches events & Audit report against the hash chain\\
\bottomrule
\end{tabularx}
\caption{Responsibility allocation: what each party is responsible for and where
that responsibility is recorded.}
\label{tab:accountability}
\end{table}

\vspace{-0.4cm}

Accountability is answered by drawing it: named roles, a signed procedure, and a
documentation chain that an auditor can follow from action to owner. The signed
standard operating procedure and the hash-chained audit trail are the concrete
mechanisms that tell a clinician, before the trial starts, exactly who is
responsible when something is wrong.

\vspace{-0.1cm}